
\subsection{System Prompt Design}
A system prompt was developed to guide the LLM-based extraction of structured drought impact events from newspaper articles. The system prompt, together with the Pydantic schema, is sent through an API to an LLM provider. The system prompt defines the general extraction task and provides the instructions that guide how the model should populate the structured output.

The system prompt contains several components. First, it defines the available impact categories from which the model can select. Second, it provides additional field guidance describing how the individual fields in the Pydantic schema should be completed. Although some of this information is also encoded directly in the Pydantic schema, it is repeated in the system prompt to improve consistency and reduce extraction errors.

The field guidance specifies how the model should extract locations, estimate recency, and identify drought impacts. It also includes a reasoning field, where the model explains why a location and associated impacts are derived from the article text. Additionally, the prompt defines seven critical rules that constrain the extraction process. The most important rule is that only realized drought consequences explicitly supported by the article should be extracted. This prevents the model from including speculative impacts, future adaptation measures, or general drought-related statements without a reported impact.

The complete system prompt is provided in Appendix~\ref{...}.

\subsection{Extraction Schema}

The extraction schema defines the structure of the LLM output and is implemented using the Python library Pydantic. Pydantic was selected because it supports schema-based output constraints across multiple LLM API providers. The schema forces the model output into a predefined JSON structure and enables automatic validation of extracted records.

After extraction, the generated output is validated within the processing pipeline. If the response cannot be parsed according to the predefined schema, the extraction request is repeated. This ensures that all extracted records follow a consistent structure before further processing.

The initial extraction schema represented each article as a collection of independent events. An event was defined as a unique combination of a location and a sentence-level fact. Each event contained a single impact label, severity or recency value, and evidence span. This structure was later revised to better represent the multi-label nature of drought impacts.



\subsection{Drought Impact Category Selection}

The first step in designing the extraction framework was determining which drought impact variables should be extracted from newspaper articles. To guide this process, a literature review on impact-based drought forecasting was conducted. Rather than developing a new categorisation from scratch, the initial framework was based on the impact categories most commonly used in previous studies.

Although the exact definitions differ between impact based forecasting studies, most organise drought impacts into four to eight broad categories. The most common groups are Agriculture and Livestock Farming, Energy and Industry, Public Water Supply, Water Resources and Water Quality, Ecosystems, and Wildfires. Many studies further aggregate these into four broad sectors: agriculture, energy, water supply, and water resources \cite{Stagge2015Modeling, lopez2025, Bulut2026Framework}. Other studies, additionally distinguish forestry, wildfires, and human health impacts \cite{Sutanto2019Moving, sodoge2023}. Although many studies adopt the 15 categories defined by the European Drought Impact Inventory (EDII), these are often aggregated into broader groups for modelling purposes.

Rather than using these broad categories, this thesis adopts a more detailed taxonomy consisting of 20 impact categories (Table~\ref{tab:impact_categories}). Because extraction is performed using zero-shot classification, the LLM is not limited by predefined training labels. A finer taxonomy therefore allows related impacts to be extracted separately, such as distinguishing Thermal/Nuclear Cooling Constraints from Industrial Water Shortages, or Forest Dieback from Wetland Loss. These detailed categories can later be aggregated into broader groups when required for forecasting or analysis.

\begin{table}[htbp]
\centering
\caption{Drought impact categories used in the extraction framework.}
\label{tab:impact_categories}
\footnotesize
\setlength{\tabcolsep}{3pt}
\renewcommand{\arraystretch}{1.0}
\begin{tabular}{lll}
\hline
Crop Failure \& Yield Reduction &
Hydropower Reduction &
Wildfire Occurrence \\
Livestock Stress \& Mortality &
Thermal/Nuclear Cooling Constraints &
Wildfire Risk Increase \\
Irrigation Shortage &
Industrial Water Shortages &
Heat \& Air Quality Health Impacts \\
Groundwater Depletion &
Inland Waterway Disruption &
Water Supply \& Sanitation Issues \\
Reservoir \& Surface Water Shortage &
Freshwater Ecosystem Degradation &
Agricultural Economic Loss \\
Water Use Restrictions &
Forest Dieback \& Vegetation Stress &
Broader Economic Disruption \\
&& Social Impacts \\
\hline
\end{tabular}
\end{table}


\newpage
\subsection{Event Representation}

The LLM extraction framework converts unstructured newspaper articles into structured drought impact events. Each event represents a reported drought impact linked to a specific location. An overview of the extracted event structure is provided in Table~\ref{tab:feature_schema_overview}. The following subsections describe the design choices behind each component in more detail.

\begin{table}[h!]
\centering
\renewcommand{\arraystretch}{1.2}
\begin{tabular}{ll}
\toprule
\textbf{Feature} & \textbf{Representation} \\
\midrule
Location & Geographic entity mentioned in the article text \\
Impact class & One or more labels from a 20-category drought impact taxonomy \\
Severity & Ordinal impact magnitude score (1--3) \\
Recency & Time since impact relative to publication date \\
Evidence & Verbatim text span supporting the extracted impact \\
Reasoning & Explanation of the location and impact extraction \\
Event structure & One event per location with associated impact categories \\
\bottomrule
\end{tabular}
\caption{Core variables in the LLM-based drought impact extraction schema.}
\label{tab:feature_schema_overview}
\end{table}
\subsubsection{Realised versus Predicted Drought Impacts}

The extraction framework focuses on realized drought impacts rather than expected consequences or adaptation measures of an drought. The objective is to ensure that the database contains observed drought impacts rather than potential future impacts.


\subsubsection{Multi-event and Multi-location Extraction}

A single article can describe drought impacts across multiple locations and sectors. To preserve spatial detail, the framework represents each distinct location mentioned in an article as a separate event, with one or more associated impact categories.

For example, an article stating that "farmers in Zeeland experienced irrigation shortages, while forests in Brabant showed signs of drought-related vegetation stress" would result in two separate events: one event linked to Zeeland with an irrigation shortage impact, and one event linked to Brabant with a forest vegetation stress impact. This differs from many traditional extraction approaches, where an article is often assigned a single dominant impact category. The proposed framework can therefore capture multiple impacts from a single sentence. Locations are extracted exactly as they are mentioned in the original text to preserve the information provided by the article.


\subsubsection{Controlled Impact Classification}

Extracted impacts are restricted to a predefined taxonomy of 20 drought impact categories (Table~\ref{tab:impact_categories}). This improves consistency and allows detailed classification of drought impacts while maintaining the possibility of later aggregation into broader impact groups.


\subsubsection{Temporal Recency Estimation}

The LLM is instructed to estimate recency based on temporal information provided within the article itself. For example, an article may state that an impact occurred "three months ago" or "last month". These references are converted into predefined recency values, providing additional temporal information for downstream impact forecasting. The main objective is to improve the predictive signal, as newspaper articles may describe impacts after they have already occurred, creating a temporal lag between the event and its reporting.


\subsubsection{Severity Assessment}

Each impact is assigned a severity score from 1 to 3 based on the reported extent and consequences of the impact. The severity guidance is adapted from the European Drought Impact Database, where 1 represents low/minor impacts, 2 represents medium/moderate impacts, and 3 represents high/extreme impacts. This provides an additional measure of impact magnitude for forecasting applications.


\subsubsection{Evidence Extraction and Reasoning}

Each event contains an evidence quote and reasoning field to improve transparency and enable manual validation of extracted impacts. The evidence quote provides the specific text supporting the extraction, for example: "Due to prolonged drought, farmers in Brabant were unable to irrigate their crops, resulting in significant yield losses."








% For our extraction schema, we use an library called pydantic, the reason we specfiiclay hcoose dfor pydatnic, for schema based output constraines is that is supported multiple llm api providers . The idea beheind an extraction schema is that the llmn is forced to output the text into schema. If it doesn't happen the llmn checks internally if it did  do that or not. Next to it after the extracition, we check in our code if the schema was pasred correclty, if not we resend the prompt to the llmn provider. The extraction schema design and the way we extracted variable's, was changed during the thesis because we thoaught anothe schema would incrasse peformacen. First our orignal extraction schema design strucutre was We have an article, we define an event as a nunique location + an fact of sentence level, then for each event we have a single impact lable and a nsingle severity or recency label, with an eivdneance span. 

% So an event o address these limitations, we introduce a post-hoc evaluation method. We sample ten manually
% verified relevant articles and evaluate model outputs in terms of false positives and false negatives at the
% extraction level. This allows us to approximate fuzzy set membership directly from predictions without
% enforcing rigid categorical boundaries.
% In parallel, we revise the model output structure to better reflect the multi-label nature of the task.
% Instead of allowing an unbounded number of extracted tuples, we enforce a hierarchical event represen-
% tation.
% The original structure was:
% • Article → Event (location + sentence-level fact)
% • Single impact label
% • Single severity or recency label
% • Single evidence span
% This is replaced by a structured formulation















% The extraction framework converts unstructured news articles into structured drought impact events while satisfying the requirements from Section 3.2, including NUTS-2/3 spatial resolution, impact categorization, severity estimation, temporal recency, and multi-location extraction.










% %0. What are the varialble's we need to extract.

% %1. How do we extract it.

% %1. What extracation framewrok we used

% %2. What is our system prompt

% %3. What is the pydantic strucutre we used
% The extraction framework converts unstructured news articles into structured drought impact events while satisfying the requirements from Section 3.2, including NUTS-2/3 spatial resolution, impact categorization, severity estimation, temporal recency, and multi-location extraction.



% An overview of the resulting event representation is provided in Table~\ref{tab:feature_schema_overview}. The framework is operationalised through a structured JSON schema enforced via Pydantic, combined with prompt-level constraints that ensure consistent extraction of all event attributes.
% \begin{table}[h!]
% \centering
% \renewcommand{\arraystretch}{1.2}
% \begin{tabular}{ll}
% \toprule
% \textbf{Feature} & \textbf{Representation} \\
% \midrule
% Location & Named geographic entity (free text, geoparsed to NUTS-2/3) \\
% Impact class & Constrained label set (20 categories) \\
% Severity & Ordinal scale (1--3) \\
% Recency & Discrete temporal buckets (0, 1, 2--24 months) \\
% Evidence & Verbatim quote (max 25 words) \\
% Reasoning & Free-text justification of classification \\
% Event structure & One record per location--impact pair \\
% \bottomrule
% \end{tabular}
% \caption{Core variables in the LLM-based drought impact extraction schema.}
% \label{tab:feature_schema_overview}
% \end{table}

% \subsection{Realised versus Predicted Drought Impacts}

% One design objective is to ensure that forecasts are non-speculative or future-oriented. News articles often include forecasts or warnings that do not correspond to observed consequences. An IBF model has no need for such impacts. To reduce the number of false positives originating from these kinds of articles, we enforce the following instruction.

% This is enforced through instructions such as:

% \begin{quote}
% ``Only include real consequences of drought; avoid vague inferences.''
% \end{quote}

% The reasoning field requires explicit justification of whether an impact is observed:

% \begin{quote}
% ``First, identify if the text describes a realised consequence of drought.''
% \end{quote}

% This improves consistency between extracted events and actual reported impacts.

% \subsection{Multi-event and Multi-location Extraction}

% Articles frequently describe multiple impacts across different locations or sectors. Aggregating these into a single record would reduce spatial and thematic resolution.

% To preserve granularity, the model generates separate entries for each location-impact pair. This is enforced by the following rule:

% \begin{quote}
% ``Each event = one distinct location + impact pair.''
% \end{quote}

% If multiple locations or impact types are present, the model is instructed to create separate entries for each:

% \begin{quote}
% ``If the article mentions multiple locations or multiple types of impacts, create a separate entry for each.''
% \end{quote}

% This enables multi-impact extraction and ensures that a single article can produce multiple outputs, which differs from traditional drought impact literature.

% \subsection{Controlled Impact Classification}

% Impact labels are constrained to predefined classes (see Appendix for the full list with prompt; up to 20 classes), covering agricultural, hydrological, ecological, infrastructural, economic, and public health impacts. The main issue with too few classes, and thus overly broad categories, is the higher likelihood of class overlap, where a single impact could belong to multiple classes, leading to ambiguity in model outputs. By starting with a higher number of classes, we reduce ambiguity and enforce more precise classification.

% The model is instructed to select only from the allowed set:

% \begin{quote}
% ``Select the most relevant category from the allowed list that defines the impact.''
% \end{quote}

% This constraint is enforced in the structured output schema using literal-type restrictions, ensuring consistently across all extracted events.

% \subsection{Temporal Recency Estimation}

% In each article, there is a referenced publication date extracted during the preprocessing step. We prompt the LLM to estimate how recent the reported impact is. In most news articles, events are described relative to the publication time, for example: crops may have been failing for two months already. This information is useful for the IBF modelling step, as it provides insight into temporal lag effects.

% We use the following mapping:

% \begin{quote}
% ``Use 0 for current/ongoing, 1 for within last month, 2--12 for months ago, or larger values for older impacts.''
% \end{quote}

% This is encoded in Pydantic as an estimated time since impact, where 0 corresponds to "current", "ongoing", or "presently".

% \subsection{Severity Assessment}

% Each event is assigned an ordinal severity score based on contextual indicators such as spatial extent, infrastructure disruption, and socioeconomic consequences. Severity is not inferred from keywords alone but from full contextual interpretation.

% The scale is defined as:

% \begin{quote}
% ``1: Localized/Moderate, 2: Widespread/Severe, 3: Large-scale/Extreme/Cascading.''
% \end{quote}

% EDID uses a similar severity scale estimate. The goal of this scale is to provide a scalar measure that allows comparison across reports. The IBF model can learn, for example, the difference between severe and less severe impacts. The idea is to compare intensity across reports and provide the IBF model with additional information for prediction. The underlying assumption is that more severe impacts are generally more strongly correlated with extreme drought events.

% \subsection{Evidence Extraction and Reasoning}

% To support interpretability, each event includes a reasoning field and a verbatim evidence quote. The evidence must directly originate from the article text and include both the impact and the location.

% This is enforced through:

% \begin{quote}
% ``Evidence must be a verbatim quote containing the location AND the impact.''
% \end{quote}

% The reasoning field explains why the extracted event qualifies as a drought impact and justifies the assigned classification and severity. Together, these fields improve transparency and enable manual validation of model outputs. When a user is validating the data or the model, it is easier to inspect the provided evidence than to read the entire article again, which can be time-consuming.
